\frfilename{mt22.tex}
\versiondate{24.2.14}
\copyrightdate{1995}

\def\chaptername{Teorema Dasar Kalkulus}
\def\sectionname{Pendahuluan}

\newchapter{22}

Dalam bab ini saya membahas salah satu sifat terpenting integral
Lebesgue. Diberikan suatu fungsi terintegralkan $f:[a,b]\to\Bbb R$, kita
dapat membentuk integral tak tentunya $F(x)=\int_a^xf(t)dt$ untuk
$x\in[a,b]$. Dua pertanyaan segera muncul. (i) Dapatkah kita mengharapkan
turunan $F'$ dari $F$ sama dengan $f$? (ii) Dapatkah kita mengenali
fungsi $F$ mana yang muncul sebagai integral tak tentu? Jawaban
yang cukup memuaskan dapat ditemukan untuk kedua pertanyaan ini: $F'=f$
hampir di mana-mana (222E) dan integral-integral tak tentu adalah
fungsi-fungsi kontinu mutlak (225E). Dalam menangani kedua pertanyaan
tersebut, kita perlu mengembangkan berbagai teknik yang menghasilkan
banyak hasil mengagumkan, baik dalam teori ukuran Lebesgue maupun dalam
topik-topik analisis real lain yang tampaknya tidak berkaitan.

Langkah pertama adalah `Teorema Vitali' (\S221), suatu argumen yang
luar biasa -- argumen ini lebih merupakan metode daripada teorema --
yang memanfaatkan sifat geometris garis bilangan real untuk mengambil
subkeluarga-subkeluarga saling lepas dari koleksi interval. Teorema ini
merupakan landasan utama bukan hanya bagi hasil-hasil dalam \S222,
melainkan juga bagi seluruh teori ukuran geometrik, yakni teori ukuran
pada ruang-ruang yang memiliki struktur geometris. Di sini saya
menggunakannya untuk menunjukkan bahwa fungsi-fungsi monoton dapat
didiferensialkan hampir di mana-mana (222A). Setelah itu, Lema Fatou dan
Teorema Konvergensi Terdominasi Lebesgue sudah cukup untuk menunjukkan
bahwa turunan suatu integral tak tentu hampir di mana-mana sama dengan
integrannya. Kita mendapati bahwa beberapa manipulasi yang tampak tidak
berbahaya terhadap fakta ini ternyata membawa kita sangat jauh; saya
menyajikannya dalam \S223.

Saya mengawali paruh kedua bab ini dengan pembahasan mengenai fungsi
`bervariasi terbatas', yakni fungsi yang dapat dinyatakan sebagai
selisih dua fungsi monoton berbatas (\S224). Bagian ini termasuk salah
satu bagian yang paling sedikit melibatkan teori ukuran dalam jilid ini;
hanya dalam 224I dan 224J ukuran dan integrasi bahkan disebut. Namun,
materi ini diperlukan untuk Bab 28 maupun untuk bagian berikutnya, dan
juga merupakan salah satu topik dasar analisis real abad kedua puluh.
\S225 membahas karakterisasi integral-integral tak tentu sebagai
fungsi-fungsi `kontinu mutlak'. Sebenarnya hal ini sekarang cukup mudah;
menggunakan kembali Teorema Vitali cukup membantu, tetapi selebihnya merupakan
penerapan langsung metode-metode yang telah digunakan sebelumnya. Paruh
kedua bagian tersebut memperkenalkan beberapa gagasan baru dalam upaya
memberikan intuisi yang lebih mendalam mengenai hakikat fungsi
kontinu mutlak. \S226 kembali kepada fungsi bervariasi terbatas dan
penguraiannya menjadi bagian `saltus', `kontinu mutlak', dan `singular';
dua bagian pertama relatif mudah ditangani, sedangkan yang terakhir agak
menyerupai fungsi Cantor.



\discrpage
